\section{Methodology (Optional)}
\label{section:algorithm}
Our method contains four components. First, we use a fixed retriever-reader backbone, such as BM25 or a small dense retriever together with a fixed reader, so that the main experimental variable is the rewrite policy itself. Second, instead of performing unrestricted free-form generation, we define a small action space for selective rewriting: keep the original query, expand entities or keywords, add temporal clarification, disambiguate a term, compress the question into keyword style, or decompose the question into two sub-queries. Third, we design a cost-aware reward that combines answer quality and retrieval quality with penalties on rewrite length and retrieved context length, encouraging the policy to rewrite only when the gain justifies the cost. Finally, for feasibility, we start from a lightweight contextual-bandit or simplified policy-gradient formulation rather than sequence-level PPO on a large language model. This design keeps the project analyzable, reproducible, and suitable for limited hardware while preserving the core RL decision problem.
